\documentclass{report}
\usepackage{amsmath,amssymb}
\setlength{\parindent}{0mm}
\setlength{\parskip}{1em}
\begin{document}
\begin{center}
\rule{6in}{1pt} \
{\large Ray Tuminaro \\
{\bf A Hybrid Operator Dependent MG / AMG Approach: Application to Ice Sheet Modeling }}

Sandia National Labs \\ MS 9159 \\ PO Box 969 \\ Livermore \\ CA 94551
\\
{\tt rstumin@sandia.gov}\\
Irina Kalashnikova\\
Mauro Perego\\
Andy Salinger\end{center}

A multigrid method is proposed that combines ideas of operator dependent
multigrid for structured grids and algebraic multigrid for unstructured
grids. It is intend for problems where the three dimensional mesh can be
viewed as an extrusion of a two dimensional unstructured mesh in a third
dimension. That is, the mesh is logically regular in the third dimension.
Our motivation comes generally from the modeling of thin structures via
finite elements and more specifically by land-based ice sheet
simulations. Extruded meshes are relatively common for thin structures
such as ice sheets and often give rise to problems with an anisotropic
nature when the thin direction mesh spacing is much smaller than the wide
direction mesh spacing. Within our approach, the first few multigrid
hierarchy levels are obtained by applying operator dependent multigrid to
semi-coarsen in a structured fashion. It is well-known that
semi-coarsening can successfully address anisotropic problems. After
sufficient structured coarsening, the resulting lower resolution mesh
contains only a single layer corresponding to a two dimensional
unstructured mesh of the wide dimensions. AMG can now be employed
successfully to create further coarse levels, as the anisotropic
phenomena is no longer present in the single layer problem.

Numerical comparisons with standard smoothed aggregation AMG are
presented demonstrating the effectiveness of the approach on several ice
sheet problems. Overall, the convergence of the new procedure is
relatively insensitive to a number of problem and mesh features.


\end{document}
